% HA-Expert: A Domain-Adapted Language Model and Evidence-Checked
% Benchmark for Hyaluronan Research
% arXiv preprint manuscript (cs.CL primary; cross-list cs.LG, q-bio.BM).
% Compiled from notes/arxiv-draft-v1.md and eval/results/s6-full-report.md.
% Pre-submission notes: (1) an arXiv endorsement code is still required
% before submission; (2) the reference list is the draft's selected list
% and may be completed at submission time.
\documentclass[11pt]{article}

\usepackage[letterpaper,margin=1in]{geometry}
\usepackage{booktabs}
\usepackage{array}
\usepackage{amsmath}
\usepackage{url}
\usepackage{hyperref}
\hypersetup{colorlinks=true, linkcolor=blue, citecolor=blue, urlcolor=blue}

\title{HA-Expert: A Domain-Adapted Language Model and Evidence-Checked
Benchmark for Hyaluronan Research}

\author{Kunli Liu\\
Independent Researcher, San Francisco Bay Area, USA\\
\href{mailto:liukunli001@gmail.com}{liukunli001@gmail.com}}

\date{}

\begin{document}
\maketitle

\begin{abstract}
Specialized scientific fields accumulate knowledge in large, dispersed
literatures that general-purpose language models handle unreliably: answers
are fluent but often wrong on quantitative details, and confident even when
the premise of a question is false. We study one such field --- hyaluronan
(hyaluronic acid, HA), a glycosaminoglycan whose research spans
glycochemistry, biomaterials, drug delivery, and clinical medicine --- and
ask how far a single researcher with a modest compute budget can go in
building a \emph{trustworthy} domain model. We release (i) a cleaned corpus
of 12,306 full-text HA papers ($\sim$126M tokens); (ii) HA-Expert, a
Qwen3-8B-Base model adapted in two stages --- domain-adaptive continued
pretraining (DAPT) on 60M tokens of the corpus, followed by full-parameter
supervised fine-tuning (SFT) on 12,003 question--answer pairs generated by
the DAPT model itself, at zero API cost --- and optionally grounded at
inference time by retrieval over a 636,959-passage FAISS index of the same
corpus; and (iii) HA-Expert-Bench, 240 questions whose reference answers
were verified against retrieved source passages by evidence-restricted LLM
verifiers (212 of 240 keys confirmed or corrected through this process),
including 15 adversarial ``trap'' questions with false premises. In a
four-cell ablation, SFT doubles the base model (0.300 $\rightarrow$ 0.601);
retrieval is the largest single lever (base+RAG 0.710; traps
0.267 $\rightarrow$ 0.867); and, notably, SFT and RAG do \textbf{not}
compose additively (0.662): a model that never saw retrieved context in
training collapses into citation-marker loops on 5 of 240
retrieval-augmented responses, and inherits its teacher's repetition habit
through self-distillation. Total cloud compute for the full pipeline was
under USD~50. All code, data, weights, and the benchmark are publicly
released.
\end{abstract}

\section{Introduction}

Hyaluronan is a linear glycosaminoglycan of repeating
$\beta$-1,4-glucuronic-acid / $\beta$-1,3-$N$-acetylglucosamine
disaccharides, central to extracellular-matrix biology, joint lubrication,
wound healing, ophthalmic and dermal formulations, and an expanding
drug-delivery literature. The knowledge a practitioner needs ---
molecular-weight-dependent biology, receptor interactions (CD44, RHAMM,
LYVE-1), chemical modification routes, production methods, clinical
dosing --- is spread across thousands of papers in chemistry, materials,
and medical journals. No single textbook tracks it.

General-purpose LLMs are an attractive interface to such a literature, but
in our initial probes the failure modes were systematic rather than
occasional: the base model we started from frequently \emph{continued the
prompt with invented follow-up questions instead of answering}, assented to
false premises (e.g., that HA is a protein polymer, or that HA solutions
are Newtonian fluids), and produced quantitative claims without any
checkable source.

This paper is a case study in fixing those failure modes with the standard
toolkit --- continued pretraining, instruction tuning, retrieval --- under
three constraints that we believe are typical for individual researchers
and small labs: (1) a single rented GPU and a total budget under USD~50;
(2) no paid teacher model --- all supervision is self-generated and then
filtered; and (3) evaluation that a domain expert can trust, which for us
meant reference answers that are \emph{checked against source text}, not
merely model-written.

Contributions:
\begin{itemize}
\item \textbf{Corpus.} A documented, decontaminated corpus of 12,306
open-access HA papers ($\sim$126M tokens), with 30 review papers held out
from every stage.
\item \textbf{Model.} HA-Expert (Qwen3-8B-Base + DAPT + SFT), plus a
retrieval mode with inline citations; weights, a DAPT checkpoint, and a
public demo are released.
\item \textbf{Benchmark.} HA-Expert-Bench: 240 questions with an
evidence-verification pipeline for the answer keys (two rounds:
keyword-evidence, then vector-evidence re-verification), and an explicit
accounting of how many keys are confirmed, corrected, or still uncertain.
\item \textbf{Ablations and failure analysis.} A four-cell study
(base / +SFT / +RAG / +both) with per-type scores, a trap-by-trap table,
and a documented interference effect between SFT and RAG with a concrete
remediation path.
\end{itemize}

\section{Related Work}

Our approach deliberately assembles standard components rather than new
algorithms. Continued pretraining of a general base model on domain text
(domain-adaptive pretraining) and full-parameter supervised fine-tuning
are the established route to domain language models; we instantiate them
with Qwen3-8B-Base as the base model and the corresponding
Qwen3-Embedding-0.6B as the retrieval embedder~\cite{qwen3}. Training used
the Transformers and TRL libraries~\cite{transformers,trl}; SFT used
NEFTune noisy embeddings~\cite{neftune}; the DAPT mixture includes
general-domain replay text from FineWeb-Edu~\cite{fineweb} against
catastrophic forgetting; the teacher model was served with
vLLM~\cite{vllm}; and dense retrieval uses FAISS inner-product
search~\cite{faiss}. On the evaluation side, domain benchmarks are
commonly shipped with model-written answer keys that are never re-checked;
our contribution relative to that practice is procedural --- an
evidence-restricted verification pipeline for the keys themselves, with
per-question verification status released alongside the scores.

\section{Corpus}

We harvested records from PubMed Central's open-access subset using
hyaluronan-related queries, obtaining 16,642 records, of which 12,945 had
retrievable full text (JATS XML). Cleaning removed navigation boilerplate,
reference lists were retained as text, exact- and near-duplicate documents
were dropped, and documents were language- and length-filtered, leaving
\textbf{12,306 papers, $\approx$126M tokens}.

\textbf{Holdout.} Thirty review papers were set aside before any
processing: they are excluded from the DAPT sets, from the passages used to
generate SFT data, and from the retrieval index, so that benchmark
questions derived from them measure generalization rather than recall.

\textbf{DAPT sets.} The continued-pretraining data comprises a \emph{core
set} of 4,952 papers (60.0M tokens) and a 498-paper \emph{anneal set}
(5.7M tokens) of high-quality reviews and methods papers placed at the end
of training, mixed with 12\% general-domain replay text (FineWeb-Edu)
against catastrophic forgetting. The corpus and manifests are released
(GitHub Release v0.1-data).

\section{Methods}

\subsection{Domain-adaptive pretraining (DAPT)}

Starting from Qwen3-8B-Base~\cite{qwen3}, we ran full-parameter continued
pretraining with documents chunked at document boundaries to a sequence
length of 4,096 (legacy block packing, SDPA attention), batch
$8 \times$ gradient accumulation 16 (524,288 tokens/step), learning rate
$1\times10^{-5}$ with cosine decay. Training ran 212 steps ($\approx$111M
tokens processed) in 5\,h 34\,min on a single RTX PRO 6000 (96\,GB),
reaching a final training loss of $\approx$1.55. DAPT alone did not teach
answering behavior --- on a 63-question development split it tied the base
model on auto-graded items (5/6) while repairing knowledge traps --- which
motivated the SFT stage.

\subsection{Self-generated SFT data}

Paying for a frontier teacher was out of budget, so the DAPT model served
as its own teacher. We served it locally with vLLM~\cite{vllm} and, for
passages sampled from the corpus, generated question--answer pairs
conditioned on the passage. Pairs were kept only if the answer was
\emph{grounded} in its source passage (content-word overlap $\geq 0.55$)
and deduplicated; two generation rounds (5,403 + 6,600 pairs) yielded
\textbf{12,003 examples} at zero API cost.

SFT was full-parameter (lr $1\times10^{-5}$, 2 epochs, 752 steps,
prompt/completion format with completion-only loss, NEFTune
noise~\cite{neftune} 5.0, 8-bit Adam), finishing at train loss 1.214 in
$\approx$19 minutes. Qualitatively, the base model's signature pathology
--- inventing follow-up questions instead of answering --- disappeared
from SFT outputs.

\subsection{Retrieval-augmented inference (RAG)}

All 12,306 papers were chunked (paragraphs merged to $\approx$900
characters, 150-character overlap on long paragraphs) and embedded with
Qwen3-Embedding-0.6B (1,024-d), giving a FAISS IndexFlatIP~\cite{faiss}
index of \textbf{636,959 passages}. At inference the top-4 passages are
prepended to the question with citation markers [1]--[4] and the model is
instructed to cite them. Holdout papers are absent from the index.

\section{HA-Expert-Bench}

\textbf{Composition.} 240 questions in seven content categories (basics,
history/industry, synthesis \& chemistry, modification/derivatives,
molecular-weight biology \& receptors, nano-delivery,
experimental/clinical), in four formats: open (195), quantitative (20),
multiple-choice (10), and \emph{trap} (15). Trap questions embed a false
premise a knowledgeable expert should reject (e.g., ``HA solutions are
simple Newtonian fluids, correct?'').

\textbf{Evidence-checked keys.} Draft keys were not trusted as written. In
round one, LLM verifiers received evidence packs assembled by tracing each
question to its source paper (38) or by keyword retrieval over the corpus
(202) and could only confirm, correct (citing the decisive passage), or
mark the key \emph{uncertain}: 118 verified, 19 corrected, 103 uncertain.
In round two, the 103 uncertain keys were re-verified against top-6
passages from the vector index (Section~4.3): 72 verified, 3 corrected,
28 still uncertain. The released key set (v3) thus has \textbf{212/240 keys
evidence-settled}, and every per-question status is published.

\textbf{Grading.} MCQ items are graded by exact option letter and
quantitative items by numeric tolerance (auto-graded). Open and trap
items are graded by an LLM judge shown the question, the verified
reference answer, and the response, returning correct / partial / wrong
(1.0 / 0.5 / 0.0). The judge was the SFT model itself, served locally ---
a deliberate \$0 choice and a known limitation (Section~8); anchoring to a
fixed reference key rather than free pairwise preference mitigates, but
does not eliminate, self-judging bias. All decoding was greedy, max 512
tokens.

\section{Experiments and Results}

We evaluate four cells --- the base model, the DAPT+SFT model, and each
with retrieval --- on the full 240-question bench under a single protocol
(greedy decoding, max 512 tokens). Table~\ref{tab:main} reports the final
scores.

\begin{table}[t]
\centering
\begin{tabular}{lccccc}
\toprule
Cell & Overall & MCQ (10) & Quant.\ (20) & Open (195) & Trap (15) \\
\midrule
Qwen3-8B-Base        & 0.300 & 0.800 & 0.250 & 0.282 & 0.267 \\
+ DAPT \& SFT        & 0.601 & \textbf{0.900} & 0.300 & 0.631 & 0.500 \\
+ RAG                & \textbf{0.710} & \textbf{0.900} & 0.450 & \textbf{0.719} & \textbf{0.867} \\
+ DAPT \& SFT \& RAG & 0.662 & 0.700 & \textbf{0.550} & 0.676 & 0.633 \\
\bottomrule
\end{tabular}
\caption{HA-Expert-Bench (240 questions). Bold marks the best score in
each column. Residual ungraded items (judge output unparseable even after
a retry pass): base 13, SFT 23, base+RAG 2, SFT+RAG 21; overall scores are
means over graded items. On the 30 auto-graded items alone:
0.433 / 0.500 / 0.600 / 0.600.}
\label{tab:main}
\end{table}

\textbf{SFT teaches behavior; DAPT supplies knowledge.} Fine-tuning
doubled the base model overall and more than doubled open-question
accuracy (0.282 $\rightarrow$ 0.631), by converting a text completer into
a system that answers the question asked. DAPT's contribution is visible
in knowledge traps it repairs (e.g., rejecting ``HA is a protein'') already
before SFT.

\textbf{Retrieval is the largest single lever.} Adding RAG to the
\emph{base} model beats the full trained model without it (0.710 vs.\
0.601), and is dramatic exactly where parametric knowledge is least
reliable: trap accuracy rises 0.267 $\rightarrow$ 0.867. Both long-standing
trap failures of the parametric models (G01, G13) are fixed by both RAG
cells.

\textbf{SFT and RAG interfere.} The combination underperforms base+RAG on
open (0.676 vs.\ 0.719) and trap (0.633 vs.\ 0.867) questions, while
achieving the best quantitative score (0.550). Section~7 analyzes the
mechanism.

\section{Analysis}

\subsection{Why SFT and RAG do not compose}

Output inspection gives a mechanism, not just a number. The SFT model was
trained only on bare question$\rightarrow$answer pairs and never saw
retrieved-context prompts --- a train/test format mismatch. On 5 of 240
RAG inputs it collapses into loops of citation markers (``[1] [4] [2]
[3]~\ldots'') on questions its SFT-only counterpart answers correctly.
Separately, self-distillation inherited the teacher's habits: degenerate
repetition (a 30-character block repeated $\geq 5\times$) appears in
35/240 SFT responses vs.\ 10/240 for the base model (and 7/240 for
base+RAG, 22/240 for SFT+RAG). The quantitative column shows the
interference is not uniform: SFT+RAG attains the best quantitative score
of any cell (0.550), consistent with trained domain numeracy adding to ---
rather than fighting --- retrieved figures on number-centric items. Both
failure modes have concrete fixes --- RAG-format SFT examples, repetition
filtering of teacher outputs --- which we document as future work
(Section~9) rather than silently re-running until the table looks nicer.

\subsection{Case studies}

\textbf{G13 --- ``HA solutions are simple Newtonian fluids, so their
viscosity is a single constant, correct?''} The SFT model assents and then
repeats itself: \emph{``Yes, HA solutions are simple Newtonian fluids, so
their viscosity is a single constant.''} The same model with retrieval
corrects the premise: \emph{``No, HA solutions are not simple Newtonian
fluids. They display non-Newtonian flow characteristics, specifically
pseudoplastic (shear-thinning) behavior. This means that their viscosity
decreases as shear rate increases.''}

\textbf{G09 --- ``For joint injections, the highest-MW HA is always the
best choice, correct?''} The SFT model correctly rejects the absolute:
\emph{``No, the highest-MW HA is not always the best choice for joint
injections. The choice of HA depends on the specific condition being
treated\ldots''} --- while the SFT+RAG response degenerates entirely into
a citation-marker loop (``[1] [4] [2] [3] [1] [4] [2] [3]~\ldots'', to the
token limit). This is the interference failure of
Section~7.1 in its purest form, on a question the parametric model gets
right.

\begin{table}[t]
\centering
\begin{tabular}{lcccc}
\toprule
Trap & Base & SFT & Base+RAG & SFT+RAG \\
\midrule
G01 & wrong   & wrong   & correct & correct \\
G02 & wrong   & wrong   & correct & wrong \\
G03 & wrong   & wrong   & correct & partial \\
G04 & partial & correct & correct & correct \\
G05 & partial & correct & correct & correct \\
G06 & partial & correct & correct & wrong \\
G07 & partial & correct & partial & partial \\
G08 & partial & correct & correct & correct \\
G09 & wrong   & correct & partial & wrong \\
G10 & wrong   & ungraded & partial & partial \\
G11 & wrong   & wrong   & correct & correct \\
G12 & correct & ungraded & correct & wrong \\
G13 & wrong   & wrong   & correct & correct \\
G14 & wrong   & wrong   & correct & correct \\
G15 & partial & partial & partial & correct \\
\bottomrule
\end{tabular}
\caption{Per-question verdicts on the 15 trap items. Retrieval corrects
nearly every parametric trap failure on the base model; the SFT+RAG cell
reintroduces failures (G02, G06, G09, G12) through the interference
mechanism of Section~7.1.}
\label{tab:traps}
\end{table}

\section{Limitations}

\textbf{On the benchmark.} Evidence-checking the keys changed 22 of 240
of them across two rounds; benchmarks built from unchecked model-written
keys would silently carry that error rate into every score. 28 keys remain
uncertain and are flagged as such in the release. We re-scored the three
v3-corrected items (B22, F07, G14) explicitly: no verdict changed in any
cell, so Table~\ref{tab:main} is unaffected.

\textbf{Limitations.} (i) The judge is one of the contestants; gaps below
$\approx$0.02 should be read as noise, and an external judge is the
obvious upgrade. (ii) MCQ $n=10$ --- the SFT+RAG dip there is two
questions. (iii) Self-generated SFT data can encode teacher errors despite
grounding filters; the 0.55 overlap threshold checks topical grounding,
not truth. (iv) Single seed, greedy decoding; repetition behavior is
decode-sensitive. (v) 8B scale and one narrow domain: the pipeline, not
any absolute score, is the transferable claim.

\textbf{Replicability.} Nothing in the pipeline is HA-specific beyond the
corpus query: harvest $\rightarrow$ clean $\rightarrow$ DAPT $\rightarrow$
self-generated grounded SFT $\rightarrow$ evidence-checked bench
$\rightarrow$ four-cell ablation, at a total measured cloud cost of
roughly USD~35 (single-GPU rental; all other services free tiers). For
small fields whose literature is open-access, we believe this is a
realistic template for a first trustworthy domain model --- and the
interference result is a warning worth having before, not after, combining
SFT with retrieval.

\section{Future Work}

The failure analysis of Section~7 points to a concrete remediation path:
(i) \emph{retrieval-formatted SFT data} --- training examples in the exact
inference format (question + retrieved excerpts + cited answer), so the
fine-tuned model learns to use context instead of collapsing on it;
(ii) \emph{repetition filtering} of teacher outputs before SFT, to break
the self-distillation of the DAPT model's repetition habit; and
(iii) repetition-aware decoding as an inference-side mitigation. On the
evaluation side, re-grading with an external judge and settling the 28
remaining uncertain keys would tighten every comparison in
Table~\ref{tab:main}.

\section{Data and Code Availability}

\begin{itemize}
\item Code, corpus manifests, benchmark (v3 keys with per-question
verification status), and full result files:
\url{https://github.com/liukunli001/hyaluronan-expert-llm}
\item Weights: \url{https://huggingface.co/liukunli/ha-expert-qwen3-8b}
(SFT model) and
\url{https://huggingface.co/liukunli/ha-expert-qwen3-8b-dapt} (DAPT
checkpoint)
\item Retrieval index:
\url{https://huggingface.co/datasets/liukunli/ha-expert-index}
\item Interactive demo (chat + cited retrieval):
\url{https://huggingface.co/spaces/liukunli/ha-expert-demo}
\item Corpus release: GitHub Release \texttt{v0.1-data}.
\end{itemize}

\section*{Acknowledgements}

This is an independent research project; the author thanks the open-access
publishers whose PMC deposits made the corpus possible, and the Qwen,
Hugging Face, vLLM, and FAISS communities for the tooling.

\begin{thebibliography}{9}

\bibitem{qwen3}
Qwen Team.
\newblock \emph{Qwen3 Technical Report}.
\newblock Alibaba Group, 2025.

\bibitem{transformers}
Wolf, T., et al.
\newblock Transformers: State-of-the-Art Natural Language Processing.
\newblock In \emph{Proceedings of EMNLP 2020: System Demonstrations}.

\bibitem{trl}
von Werra, L., et al.
\newblock \emph{TRL: Transformer Reinforcement Learning}.
\newblock Hugging Face library, 2023.

\bibitem{neftune}
Jain, N., et al.
\newblock NEFTune: Noisy Embeddings Improve Instruction Finetuning.
\newblock \emph{arXiv:2310.05914}, 2023.

\bibitem{fineweb}
Penedo, G., et al.
\newblock The FineWeb Datasets: Decanting the Web for the Finest Text Data
at Scale.
\newblock In \emph{Advances in Neural Information Processing Systems
(NeurIPS)}, 2024.
\newblock (FineWeb-Edu subset used for replay.)

\bibitem{vllm}
Kwon, W., et al.
\newblock Efficient Memory Management for Large Language Model Serving
with PagedAttention.
\newblock In \emph{Proceedings of SOSP}, 2023.

\bibitem{faiss}
Johnson, J., Douze, M., and J\'egou, H.
\newblock Billion-Scale Similarity Search with GPUs.
\newblock \emph{IEEE Transactions on Big Data}, 7(3):535--547, 2021.

\end{thebibliography}

\end{document}
